% =====================================================================
%  Section 6: Conclusion
% =====================================================================
\section{Conclusion}
\label{sec:conclusion}

We presented \textsc{TADS} (Trajectory-Anchored Dynamic Selection),
whose central contribution is the \emph{trajectory anchor}: a
per-layer capability direction extracted online from the model's
own hidden-state trajectory at each refresh step, with no
in-domain seed, labels, validation set, or auxiliary scorer
required. The anchor is what makes \textsc{TADS} different from
prior dynamic IT-data selectors -- it captures direction in
representation space rather than only the magnitude of a scalar
composite reward. The remaining components of the selection rule
(a deterministic pool-normalised calibrated utility
$\tilde{\mathcal{R}}_i^{(t)} \in (0,1)$, a single multiplicative
combination, and one hyperparameter $\lambda$) exist to integrate
the anchor with the existing composite-reward signal in a
reproducible way, with $\lambda{=}0$ disabling the anchor and
recovering a clean base ranking.

We proved a stability guarantee (Theorem~\ref{thm:stability}):
under a positive eigenvalue gap, the trajectory anchor changes by
at most a learning-rate-scaled amount per epoch, and this bound
transfers to the ranking criterion. To our knowledge, this is the
first stability guarantee for a dynamic IT data-selection method,
and it is base-agnostic rather than tied to a specific utility
transform.

On the LLaMA-2-7B + Alpaca-GPT4 setting at $10\%$ selection,
\textsc{TADS} improves over the calibrated-utility base on the
completed metrics (MMLU, GSM8K, SVAMP, MBPP, TydiQA, BBH) while
adding under $0.5\%$ wall-clock overhead. A broader study --
across model scales from 0.5B to 14B, multiple instruction-tuning
datasets, and an eight-task benchmark suite -- is in progress and
described in App.~\ref{app:inprogress}; it is designed to test
whether the gains hold consistently across models, datasets, and
evaluation regimes under a single \textsc{TADS} configuration.
